% Options for packages loaded elsewhere
\PassOptionsToPackage{unicode}{hyperref}
\PassOptionsToPackage{hyphens}{url}
\PassOptionsToPackage{dvipsnames,svgnames,x11names}{xcolor}
%
\documentclass[
  10pt,
]{article}
\usepackage{amsmath,amssymb}
\usepackage{iftex}
\ifPDFTeX
  \usepackage[T1]{fontenc}
  \usepackage[utf8]{inputenc}
  \usepackage{textcomp} % provide euro and other symbols
\else % if luatex or xetex
  \usepackage{unicode-math} % this also loads fontspec
  \defaultfontfeatures{Scale=MatchLowercase}
  \defaultfontfeatures[\rmfamily]{Ligatures=TeX,Scale=1}
\fi
\usepackage{lmodern}
\ifPDFTeX\else
  % xetex/luatex font selection
  \setmainfont[]{DejaVu Serif}
  \setsansfont[]{DejaVu Sans}
  \setmonofont[]{DejaVu Sans Mono}
  \setmathfont[]{Latin Modern Math}
\fi
% Use upquote if available, for straight quotes in verbatim environments
\IfFileExists{upquote.sty}{\usepackage{upquote}}{}
\IfFileExists{microtype.sty}{% use microtype if available
  \usepackage[]{microtype}
  \UseMicrotypeSet[protrusion]{basicmath} % disable protrusion for tt fonts
}{}
\makeatletter
\@ifundefined{KOMAClassName}{% if non-KOMA class
  \IfFileExists{parskip.sty}{%
    \usepackage{parskip}
  }{% else
    \setlength{\parindent}{0pt}
    \setlength{\parskip}{6pt plus 2pt minus 1pt}}
}{% if KOMA class
  \KOMAoptions{parskip=half}}
\makeatother
\usepackage{xcolor}
\usepackage[margin=2.0cm]{geometry}
\usepackage{longtable,booktabs,array}
\usepackage{calc} % for calculating minipage widths
% Correct order of tables after \paragraph or \subparagraph
\usepackage{etoolbox}
\makeatletter
\patchcmd\longtable{\par}{\if@noskipsec\mbox{}\fi\par}{}{}
\makeatother
% Allow footnotes in longtable head/foot
\IfFileExists{footnotehyper.sty}{\usepackage{footnotehyper}}{\usepackage{footnote}}
\makesavenoteenv{longtable}
\setlength{\emergencystretch}{3em} % prevent overfull lines
\providecommand{\tightlist}{%
  \setlength{\itemsep}{0pt}\setlength{\parskip}{0pt}}
\setcounter{secnumdepth}{-\maxdimen} % remove section numbering
% doc_header.tex -- LaTeX preamble additions used by build_docs.sh for every
% Markdown -> LaTeX -> PDF conversion (supplements, README, submission files).
\usepackage{etoolbox}
% longtable may not be loaded by pandoc when a document has no table
\usepackage{longtable}
% bold vectors: \bm -> \symbf under unicode-math (XeLaTeX), bm package otherwise
\makeatletter
\@ifpackageloaded{unicode-math}{\providecommand{\bm}[1]{\symbf{#1}}}{\usepackage{bm}}
\makeatother
% smaller type in long tables so that wide audit tables fit the text width;
% inside a table a line break is allowed after an underscore of plain text
% (pandoc writes it as \_), so that long identifiers wrap inside their column
% rows are set a little apart, so that symbols with indices do not touch
\AtBeginEnvironment{longtable}{\footnotesize\renewcommand{\arraystretch}{1.45}%
  \let\sptunderscore\_\def\_{\sptunderscore\allowbreak}}
\setlength{\LTcapwidth}{\textwidth}
% allow line breaks after slashes and hyphens inside table cells
\usepackage{xurl}
% code blocks: break over-long lines instead of running into the margin.
% Plain ``` blocks are typeset by pandoc as {verbatim}, highlighted ones
% (```bash) as {Highlighting}; both are re-declared as breakable fancyvrb
% environments.  The continuation line is marked with a small hooked arrow.
\usepackage{fvextra}
\RecustomVerbatimEnvironment{verbatim}{Verbatim}{breaklines,breakanywhere,fontsize=\small}
\makeatletter
\@ifundefined{Highlighting}{}{%
  \RecustomVerbatimEnvironment{Highlighting}{Verbatim}{breaklines,breakanywhere,fontsize=\small,commandchars=\\\{\}}}
\makeatother
% inline code (pandoc writes it as \texttt{...}): never hyphenate -- an inserted
% hyphen could be read as part of a file name -- but allow a line break after
% an underscore, slash, comma, colon or equals sign, so that long paths and
% column lists wrap inside the text width and inside table columns.
% Inside a formula (e.g. the name of a notebook symbol in an equation) the same
% command must first leave math mode, which \text of amsmath does.
\ExplSyntaxOn
\tl_new:N \l__spt_code_tl
\cs_new_protected:Npn \__spt_texttt:n #1
  {
    \tl_set:Nn \l__spt_code_tl {#1}
    \regex_replace_all:nnN { ([\/\,\:\=]) } { \1 \c{allowbreak} } \l__spt_code_tl
    \regex_replace_all:nnN { \c{_} } { \c{_} \c{allowbreak} } \l__spt_code_tl
    \group_begin: \ttfamily \hyphenchar \font = -1 \scan_stop: \l__spt_code_tl \group_end:
  }
\RenewDocumentCommand \texttt { m }
  {
    \mode_if_math:TF { \text { \__spt_texttt:n {#1} } } { \__spt_texttt:n {#1} }
  }
\ExplSyntaxOff
% A figure is written in the Markdown source as an image paragraph followed by
% a caption paragraph ("**Figure S1.** ..."), so that the caption is also
% visible on GitHub.  An image that starts a paragraph is centered.
\usepackage{graphicx}% pandoc loads it only for documents that contain an image
\NewCommandCopy\sptincludegraphics\includegraphics
\RenewDocumentCommand\includegraphics{O{} m}{%
  \ifvmode\begin{center}\sptincludegraphics[#1]{#2}\end{center}%
  \else\sptincludegraphics[#1]{#2}\fi}
% keep figures where they are referenced
\usepackage{float}
\floatplacement{figure}{H}
% running header with the document title
\usepackage{fancyhdr}
\pagestyle{fancy}
\fancyhf{}
\fancyhead[L]{\small PROJECT spacetimes --- documentation and provenance}
\fancyhead[R]{\small\thepage}
\renewcommand{\headrulewidth}{0.4pt}
% 16x16 matrices: tighter columns so that a full matrix fits the text width
\setlength{\arraycolsep}{3.2pt}
% The documents of this project contain long tagged symbols, file paths and
% identifiers that cannot be hyphenated.  A paragraph is therefore allowed to
% have looser lines before TeX gives up and lets a line stick out of the margin.
\AtBeginDocument{\tolerance=2000 \emergencystretch=5em \hbadness=2000}
\ifLuaTeX
  \usepackage{selnolig}  % disable illegal ligatures
\fi
\IfFileExists{bookmark.sty}{\usepackage{bookmark}}{\usepackage{hyperref}}
\IfFileExists{xurl.sty}{\usepackage{xurl}}{} % add URL line breaks if available
\urlstyle{same}
\hypersetup{
  pdftitle={P3 --- The Mathematica notebooks of the project folder: which were used, and how},
  pdfauthor={Prepared for Patrick Lee Nash, Ph.D.},
  colorlinks=true,
  linkcolor={blue},
  filecolor={Maroon},
  citecolor={blue},
  urlcolor={blue},
  pdfcreator={LaTeX via pandoc}}

\title{P3 --- The Mathematica notebooks of the project folder: which
were used, and how}
\usepackage{etoolbox}
\makeatletter
\providecommand{\subtitle}[1]{% add subtitle to \maketitle
  \apptocmd{\@title}{\par {\large #1 \par}}{}{}
}
\makeatother
\subtitle{Provenance required by the task specification (PROJECT
spacetimes)}
\author{Prepared for Patrick Lee Nash, Ph.D.}
\date{November 12, 2026}

\begin{document}
\maketitle

\hypertarget{what-this-document-is-for}{%
\section{1. What this document is for}\label{what-this-document-is-for}}

The task specification asks for a provenance that lists every
Mathematica notebook that was used and says how it was used. The project
folder contains fourteen notebooks. This document lists them, describes
the tools with which they were read, and states, cell by cell where it
matters, what the project took from each. The \emph{symbols} of the
notebooks, and the many places where two notebooks use the same symbol
for different objects, are the subject of the MASTER provenance
(\texttt{provenance/P1\_master\_symbols}).

The result in short:

\begin{itemize}
\tightlist
\item
  \texttt{4+4.nb} is the source of the eight real \(16\times16\) Dirac
  matrices of the 4+4 spacetime, as the task specification prescribes.
\item
  \texttt{7+1.nb} is the source of the eight real \(16\times16\) Dirac
  matrices of the 7+1 spacetime, as the task specification prescribes,
  and is quoted in a few places for context.
\item
  The meaning of the eight coordinates is taken from the head of
  \texttt{4+4.nb} and from
  \texttt{Dirac\_Einstein-Lovelock\_4+4\_Nash.nb}.
\item
  The other eleven notebooks were read for their conventions and for the
  survey of symbols. Nothing is taken from them; most belong to the
  second stage of the author's programme (gravity), which the task
  specification says is not to be implemented.
\item
  \textbf{Mathematica was not run.} The project reads what Mathematica
  \emph{printed} in the notebooks, rebuilds every matrix from its
  definition with its own program (Rust), compares the two entry by
  entry, and repeats the essential checks with an independent program
  (Python). Nothing in the project depends on evaluating a notebook
  again.
\end{itemize}

The notebooks are not part of the repository (they are the author's
unpublished work); \texttt{provenance/inputs\_sha256.txt} identifies the
copies that were read.

\hypertarget{how-the-notebooks-were-read}{%
\section{2. How the notebooks were
read}\label{how-the-notebooks-were-read}}

A notebook file (\texttt{.nb}) is a text file that contains one large
expression: a list of cells, each holding typeset ``boxes''. Three small
programs, all in the repository, read these files without Mathematica.

\begin{longtable}[]{@{}
  >{\raggedright\arraybackslash}p{(\columnwidth - 4\tabcolsep) * \real{0.2842}}
  >{\raggedright\arraybackslash}p{(\columnwidth - 4\tabcolsep) * \real{0.4421}}
  >{\raggedright\arraybackslash}p{(\columnwidth - 4\tabcolsep) * \real{0.2737}}@{}}
\toprule\noalign{}
\begin{minipage}[b]{\linewidth}\raggedright
Program
\end{minipage} & \begin{minipage}[b]{\linewidth}\raggedright
What it does
\end{minipage} & \begin{minipage}[b]{\linewidth}\raggedright
Used for
\end{minipage} \\
\midrule\noalign{}
\endhead
\bottomrule\noalign{}
\endlastfoot
\texttt{tools/nb2txt.py} & converts a notebook to plain text, cell by
cell, with a running cell number and the style of each cell (Input,
Output, Text, \ldots); typeset input appears in linear form & reading;
every citation ``cell \emph{n}'' in the project refers to its
numbering \\
\texttt{tools/extract\_nb\_matrices.py} & reads the integer matrices
that Mathematica printed in given output cells &
\texttt{data/notebook\_matrices\_4+4.txt},
\texttt{data/notebook\_matrices\_7+1.txt}, against which the programs
compare their matrices \\
\texttt{tools/scan\_notebook\_symbols.py} & lists the symbols that each
notebook defines & the survey of symbols of P1 \\
\end{longtable}

To repeat the reading (from the root of the repository, with the
notebooks in that folder):

\begin{verbatim}
python3 tools/nb2txt.py "4+4.nb" /tmp/4+4.txt 300          # text of a notebook; outputs cut after 300 characters
python3 tools/extract_nb_matrices.py "4+4.nb" 13 136 163 165 314 > data/notebook_matrices_4+4.txt
python3 tools/extract_nb_matrices.py "7+1.nb" 204 > data/notebook_matrices_7+1.txt
python3 tools/scan_notebook_symbols.py                     # symbols of all notebooks
\end{verbatim}

\hypertarget{the-fourteen-notebooks}{%
\section{3. The fourteen notebooks}\label{the-fourteen-notebooks}}

``Cells'' is the number of cells, with the number of input cells in
parentheses.

\begin{longtable}[]{@{}
  >{\raggedright\arraybackslash}p{(\columnwidth - 4\tabcolsep) * \real{0.0617}}
  >{\raggedright\arraybackslash}p{(\columnwidth - 4\tabcolsep) * \real{0.7901}}
  >{\raggedright\arraybackslash}p{(\columnwidth - 4\tabcolsep) * \real{0.1481}}@{}}
\toprule\noalign{}
\begin{minipage}[b]{\linewidth}\raggedright
Id
\end{minipage} & \begin{minipage}[b]{\linewidth}\raggedright
File
\end{minipage} & \begin{minipage}[b]{\linewidth}\raggedright
Cells
\end{minipage} \\
\midrule\noalign{}
\endhead
\bottomrule\noalign{}
\endlastfoot
N01 & \texttt{4+4.nb} & 806 (366) \\
N02 & \texttt{7+1.nb} & 1402 (189) \\
N03 & \texttt{Bessel\ product.nb} & 184 (88) \\
N04 & \texttt{Dirac\_Einstein-Lovelock\_4+4\_Nash.nb} & 403 (112) \\
N05 & \texttt{Dirac\_Nash\_Identity.nb} & 400 (192) \\
N06 & \texttt{Einstein-Lovelock-4+4-Nash.nb} & 343 (82) \\
N07 & \texttt{Einstein-Lovelock-7+1-Nash.nb} & 342 (82) \\
N08 & \texttt{Einstein-Lovelock-Eternal.nb} & 316 (64) \\
N09 & \texttt{Generalized\ \_Kronecker\_Delta\_4+4.nb} & 218 (57) \\
N10 &
\texttt{Pair\_Creation\_of\_Universes\_WaveFunctionOfUniverse-4+4-Einstein-Lovelock-Nash.nb}
& 2098 (991) \\
N11 & \texttt{WaveFunctionUniverse-4+4-Einstein-Lovelock-Nash.nb} & 929
(452) \\
N12 & \texttt{WaveFunctionUniverse-4+4-Einstein-Lovelock.nb} & 1317
(615) \\
N13 & \texttt{det-g.nb} & 20 (8) \\
N14 & \texttt{field\_equations.nb} & 271 (69) \\
\end{longtable}

What the notebooks contain, and what the project did with them:

\begin{longtable}[]{@{}
  >{\raggedright\arraybackslash}p{(\columnwidth - 4\tabcolsep) * \real{0.0641}}
  >{\raggedright\arraybackslash}p{(\columnwidth - 4\tabcolsep) * \real{0.5128}}
  >{\raggedright\arraybackslash}p{(\columnwidth - 4\tabcolsep) * \real{0.4231}}@{}}
\toprule\noalign{}
\begin{minipage}[b]{\linewidth}\raggedright
Id
\end{minipage} & \begin{minipage}[b]{\linewidth}\raggedright
What it contains
\end{minipage} & \begin{minipage}[b]{\linewidth}\raggedright
Use
\end{minipage} \\
\midrule\noalign{}
\endhead
\bottomrule\noalign{}
\endlastfoot
N01 & the real Dirac matrices of 4+4 and their \(8\times8\) and
\(4\times4\) building blocks; split octonions; complex structures &
\textbf{source of the 4+4 matrices} (Sec. 4) \\
N02 & octonions; real Dirac matrices of 7+1; Dirac equation of a
16-component field; in 27 sections also gravity, a brane picture,
cosmology & \textbf{source of the 7+1 matrices}; quoted for context
(Sec. 5) \\
N03 & identities for products of Bessel functions & not used \\
N04 & Einstein--Lovelock field equations on 4+4 with the package xAct; a
text on the origin of the work & \textbf{meaning of the coordinates}
\(x0,\dots,x7\); the lineage quoted in the Introduction of the
manuscripts \\
N05 & identities of real \(4\times4\) matrices & read; not used \\
N06 & field equations on 4+4 (xAct) & read; not used \\
N07 & field equations on 7+1 (xAct) & read; not used \\
N08 & a solution with inflating and deflating dimensions & read for the
picture of Stage 2 that the task specification describes; not used \\
N09 & generalized Kronecker delta, Lovelock tensors & read; not used \\
N10 & a wave function of the universe on 4+4; repeats the construction
of N01 & read; not used \\
N11 & as N10 & read; not used \\
N12 & as N10 & read; not used \\
N13 & determinant of a diagonal metric & read; not used \\
N14 & field equations on 4+4 (xAct) & read; not used \\
\end{longtable}

``Read'' means: the notebook was converted completely to text, its
headings and text cells were read, and every cell entered the survey of
symbols. Only for N01, N02 and N04 does a statement of the project rest
on the content of a notebook; the following sections list these
statements.

\hypertarget{what-was-taken-from-44.nb-n01}{%
\section{\texorpdfstring{4. What was taken from \texttt{4+4.nb}
(N01)}{4. What was taken from 4+4.nb (N01)}}\label{what-was-taken-from-44.nb-n01}}

The project rebuilds the matrices from the definitions in the input
cells of the notebook and compares the result with the matrices that
Mathematica printed.

\begin{longtable}[]{@{}
  >{\raggedright\arraybackslash}p{(\columnwidth - 6\tabcolsep) * \real{0.3803}}
  >{\raggedright\arraybackslash}p{(\columnwidth - 6\tabcolsep) * \real{0.2394}}
  >{\raggedright\arraybackslash}p{(\columnwidth - 6\tabcolsep) * \real{0.2394}}
  >{\raggedright\arraybackslash}p{(\columnwidth - 6\tabcolsep) * \real{0.1408}}@{}}
\toprule\noalign{}
\begin{minipage}[b]{\linewidth}\raggedright
Object of the notebook
\end{minipage} & \begin{minipage}[b]{\linewidth}\raggedright
Symbol of the project
\end{minipage} & \begin{minipage}[b]{\linewidth}\raggedright
Printed in
\end{minipage} & \begin{minipage}[b]{\linewidth}\raggedright
Compared by check
\end{minipage} \\
\midrule\noalign{}
\endhead
\bottomrule\noalign{}
\endlastfoot
\texttt{s4by4{[}h{]}}, \texttt{t4by4{[}h{]}}, \texttt{ζ{[}A{]}},
\texttt{τ{[}A{]}}, \texttt{τ} with an overbar, \texttt{σ} \(=\)
\texttt{σ88} & \(s_h\), \(t_h\), \(\tau_{1644}^A\),
\(\bar\tau_{1644}^A\), \(\sigma_{88}\) & Out{[}76{]}, Out{[}77{]}
(relations) & \texttt{44-B1} to \texttt{44-B4} \\
\texttt{T16\^{}A{[}0{]}} to \texttt{T16\^{}A{[}7{]}} &
\(\gamma_{1644}^A\) & cell 314, Out{[}211{]} & \texttt{44-N1} \\
\texttt{T16\^{}A{[}8{]}} & \(\omega_{1644}\) & cell 136, Out{[}99{]} &
\texttt{44-N2} \\
\texttt{η4488} & \(\eta_{1644}\) & cell 13, Out{[}9{]} &
\texttt{44-N3} \\
\texttt{P\_L}, \texttt{P\_R} & \(P_{1644}^{L}\), \(P_{1644}^{R}\) &
cells 163 and 165, Out{[}112{]}, Out{[}113{]}; Out{[}114{]} (relations)
& \texttt{44-P3}, \texttt{44-P2} \\
\texttt{σ16} & \(C_{1644}^{-}\) & Out{[}105{]}, Out{[}106{]} (relations)
& \texttt{44-S1}, \texttt{44-S4} \\
\texttt{JT4} \(=\) \texttt{T16\^{}A{[}4{]}} & \(\gamma_{1644}^4\) as a
complex structure & --- & \texttt{44-S5} \\
\end{longtable}

The meaning of the coordinates (\(x0\) hidden space, \(x1,x2,x3\) space,
\(x4\) time, \(x5,x6,x7\) additional times) is stated in a text cell at
the head of the notebook.

Not taken from N01: its multiplication table of the split octonions, its
frames (\texttt{E\^{}A\_a}), and its search for complex structures. The
Dirac \emph{equation} of the field dirac1644, its current, its
Lagrangian and everything that follows are not in the notebook; they are
derived in ``dirac matrices'', Part 1.

\hypertarget{what-was-taken-from-71.nb-n02}{%
\section{\texorpdfstring{5. What was taken from \texttt{7+1.nb}
(N02)}{5. What was taken from 7+1.nb (N02)}}\label{what-was-taken-from-71.nb-n02}}

The notebook labels the time by the index 0; the project labels it by 4,
as the task specification demands, and relates the two labellings by the
exchange \(\mathrm{swap}_{04}\) of the index values 0 and 4.

\begin{longtable}[]{@{}
  >{\raggedright\arraybackslash}p{(\columnwidth - 6\tabcolsep) * \real{0.3803}}
  >{\raggedright\arraybackslash}p{(\columnwidth - 6\tabcolsep) * \real{0.2394}}
  >{\raggedright\arraybackslash}p{(\columnwidth - 6\tabcolsep) * \real{0.2394}}
  >{\raggedright\arraybackslash}p{(\columnwidth - 6\tabcolsep) * \real{0.1408}}@{}}
\toprule\noalign{}
\begin{minipage}[b]{\linewidth}\raggedright
Object of the notebook
\end{minipage} & \begin{minipage}[b]{\linewidth}\raggedright
Symbol of the project
\end{minipage} & \begin{minipage}[b]{\linewidth}\raggedright
Where in the notebook
\end{minipage} & \begin{minipage}[b]{\linewidth}\raggedright
Compared by check
\end{minipage} \\
\midrule\noalign{}
\endhead
\bottomrule\noalign{}
\endlastfoot
\texttt{Γ17{[}0{]}} to \texttt{Γ17{[}7{]}} (real set) & --- & cell 204,
Out{[}178{]}; Out{[}173{]} (relations) & \texttt{71-N1},
\texttt{71-N3} \\
\texttt{Γ17up{[}A{]}} \(=\) \texttt{η17}\(^{AA}\) \texttt{Γ17{[}A{]}} &
\(\gamma_{1671}^{\mathrm{swap}_{04}(A)}\) & its Sec. 5.3 &
\texttt{71-N2} \\
\texttt{τO{[}A{]}} (octonionic \(8\times8\) matrices) &
\(\tau_{1671}^{\mathrm{swap}_{04}(A)}\) & its Sec. 3.1; Out{[}175{]}
(block form) & \texttt{71-B1}, \texttt{71-B2}, \texttt{71-D6} \\
the construction of \texttt{Γ17} from \texttt{T16} and \texttt{S44} &
relation between the two sets of the project & Out{[}176{]} &
\texttt{71-B3}, \texttt{71-B4} \\
\texttt{Γ17{[}8{]}} \(=\) \texttt{ω17}, \texttt{χ17}, \texttt{PL17},
\texttt{PR17} & \(\omega_{1671}\), \(\chi_{1671}\), \(P_{1671}^{L,R}\) &
its Secs. 4.6 and 7.10 & \texttt{71-V1}, \texttt{71-P1},
\texttt{71-P4} \\
\texttt{βD}, \texttt{αD{[}i{]}} & \(\beta_{1671}\), \(\alpha_{1671}^j\)
& its Sec. 7.1 & \texttt{71-D6} \\
\texttt{dirac16complex} and its Dirac equation & \(\psi_{1671}\)
(dirac1671), Eq. (2.9) of ``dirac matrices'' & its Sec. 7.2 &
\texttt{71-D1} and following \\
\texttt{S44} & \(C_{1644}^{+}\) & its Sec. 4.1 & \texttt{44-S2} \\
\end{longtable}

Statements of the notebook that the project quotes without using them in
a calculation:

\begin{longtable}[]{@{}
  >{\raggedright\arraybackslash}p{(\columnwidth - 4\tabcolsep) * \real{0.6528}}
  >{\raggedright\arraybackslash}p{(\columnwidth - 4\tabcolsep) * \real{0.0972}}
  >{\raggedright\arraybackslash}p{(\columnwidth - 4\tabcolsep) * \real{0.2500}}@{}}
\toprule\noalign{}
\begin{minipage}[b]{\linewidth}\raggedright
Statement
\end{minipage} & \begin{minipage}[b]{\linewidth}\raggedright
Section of \texttt{7+1.nb}
\end{minipage} & \begin{minipage}[b]{\linewidth}\raggedright
Quoted in
\end{minipage} \\
\midrule\noalign{}
\endhead
\bottomrule\noalign{}
\endlastfoot
With the signs of N01 (space-like \(=+1\)) no real \(16\times16\)
matrices exist for 7+1; with the opposite signs they do & 4.4 & ``dirac
matrices'', Sec. 2.4 (the project proves the positive part itself,
checks \texttt{71-R1}, \texttt{71-C1}, \texttt{71-C3}) \\
Plane waves, spin sums and quantization of the field were derived
symbolically & 7.4 to 7.7 & ``dirac matrices'', Sec. 2.8.2 (tested
numerically there) \\
The two coupled 8-component equations & 7.11 & ``dirac matrices'', Sec.
2.8.3 \\
A brane picture with the Standard Model on a brane & 17, 18 & Supplement
S2, as an alternative to the hypothesis examined there \\
A 4-torus and a 4-sphere as extra spaces & 25, 26 & Supplement S2, for
the isometries of possible extra spaces \\
A mass window for the unpaired field as dark matter & 24 & Supplement
S3, Sec. 4.3, for comparison \\
\end{longtable}

Not taken from N02: its complex set \texttt{Γ71}, its \(32\times32\)
set, its Kohn--Sham theory of a self-coupled field (its Secs. 8, 9, 15,
16), its gravity (Secs. 11 to 14, 20 to 23) and its cosmology. The
interaction of this project is a different one: it is mediated by the
Vorton-Scalar, which the notebook does not contain.

The notebook says of itself that it was generated by a script and
evaluated in Mathematica, and that each of its claims is tested by a
cell of the notebook. The project does not rely on those tests: every
statement about the matrices that it uses is tested again by its own
programs.

\hypertarget{what-was-taken-from-the-other-notebooks}{%
\section{6. What was taken from the other
notebooks}\label{what-was-taken-from-the-other-notebooks}}

\textbf{N04.} Its text cells state the meaning of the coordinates in the
same words as N01, and describe the origin of the real representations
(Dirac's representation of the covering group of \(SO(3,2)\), extended
by the author to \(SO(3,3)\), \(SO(4,4)\) and real \(16\times16\)
matrices); the Introduction of the manuscripts repeats this in one
sentence, with the references that P2 lists. Its field equations are not
used.

\textbf{N06 to N09, N13, N14 (gravity).} The task specification
describes Stage 2 of the author's programme (three space dimensions
inflate, three time dimensions deflate, the hidden dimension curls up
into a ring) and says: be aware of this, but do not implement it. These
notebooks are that stage. The project uses the description only where it
discusses what Stage 2 might change (Supplement S1, Secs. 6.3 and 8;
Supplement S2; the discussion sections of the manuscripts), and takes no
equation from them.

\textbf{N10 to N12.} They repeat the matrices of N01 with identical
statements (the survey of symbols shows this) and add a wave function of
the universe, which is not part of this project.

\textbf{N03, N05.} Unrelated to the task, or concerned with
four-dimensional matrices.

\hypertarget{checksums}{%
\section{7. Checksums}\label{checksums}}

The following lines of \texttt{provenance/inputs\_sha256.txt} identify
the fourteen copies (SHA-256; verify with
\texttt{shasum\ -a\ 256\ \textless{}file\textgreater{}} on macOS):

\begin{verbatim}
740023233eaa0cd24ca0e2e7010094671b1d314b698114c08db94050d3e0aed5  4+4.nb
41d252f9852f210c228f45e1ac956d2d7ef26faeb0062a5880280e1aae0effc7  7+1.nb
5e3d97c18050178f74068925c540d6ac5ca0eaa8f9416cc53b96f6a961959b81  Bessel product.nb
cdbb6a2c0b31ffd6a4fd951746a1e1c3ea5748afd371edc5095cea7e34cf79c0  Dirac_Einstein-Lovelock_4+4_Nash.nb
66267b2bf43dc711c6228822948db4da8f4b26b1a74d5839c049a35c5314e842  Dirac_Nash_Identity.nb
e749997952af0ba64e347864cab209d140cc361cd059effaadf3d68410f1f96b  Einstein-Lovelock-4+4-Nash.nb
af47a11d9a5d9f56a46a81b0613626ed3b5b7fd13111ed21bd4d3feee763acdc  Einstein-Lovelock-7+1-Nash.nb
a0e01b0101bcad8ee252591065b3999b0efe816b5eb306af2c4e0abd19c888cb  Einstein-Lovelock-Eternal.nb
23bb4e0c70943e766d9b081a3a399ef29664889aad088041329ce2e065b80afb  Generalized _Kronecker_Delta_4+4.nb
5ee5cb2a95146136ee65636a4aef174f303b6c41da130d2c57e4684f9c8ff69f  Pair_Creation_of_Universes_WaveFunctionOfUniverse-4+4-Einstein-Lovelock-Nash.nb
d06875b70d0926bd12fbb3d24cc14f2b425be7feaa9956b581319ff54004fa97  WaveFunctionUniverse-4+4-Einstein-Lovelock-Nash.nb
3e0ef380cc36198cbc05da72de26a4218c0a64ab0711fa179a1701c77ca0fb8b  WaveFunctionUniverse-4+4-Einstein-Lovelock.nb
afca8dcfd94394eb3acb0ba60444968a062ba7c86bbf8a09e6037a6f87c74e9d  det-g.nb
3bd90c1b93aa29c53e1388a4f73f6ec0645227124dbc8ee217c21fdde6d9a3f9  field_equations.nb
\end{verbatim}

\end{document}
